% =============================================================================
% Tables of lattice invariants: constructor
% =============================================================================
% A table of lattice invariants is the grid of integer points (x, y) of two
% invariants, with a vertex mark at each point where a lattice with those
% invariants occurs; the mark says which values of a third invariant occur
% there. Nikulin's table of 2-elementary lattices by (r, a, delta),
% figures/objects/nikulin/two-elementary-table.tikz, is one such object.
% \input by dzg-constructors.tex, so "@" is a letter here.
%
% Command:
%   \latticetable{<x name>}{<x max>}{<y name>}{<y max>}{<x/y/mark, ...>}
%       the grid [0, x max + 1] x [0, y max + 1] with axes named <x name> and
%       <y name> and ticks 1, ..., x max and 1, ..., y max; at each listed
%       point a node[vertex=<mark>] (mark white, black or doubled) named
%       -<x>-<y>.
%
% Key:
%   lattice table labels=axes|none   axis names and ticks (default axes)

\newif\ifdzg@lt@labels
\dzg@lt@labelstrue
\tikzset{
  lattice table labels/.is choice,
  lattice table labels/axes/.code={\dzg@lt@labelstrue},
  lattice table labels/none/.code={\dzg@lt@labelsfalse},
}

\newcommand\latticetable[5]{%
  \pgfmathtruncatemacro\dzg@lt@xgrid{#2+1}%
  \pgfmathtruncatemacro\dzg@lt@ygrid{#4+1}%
  \pgfmathsetmacro\dzg@lt@xend{#2+1.6}%
  \pgfmathsetmacro\dzg@lt@yend{#4+1.6}%
  \draw[lattice grid, step={(1,1)}] (0,0) grid (\dzg@lt@xgrid,\dzg@lt@ygrid);
  \draw[->] (0,0) -- (\dzg@lt@xend,0);
  \draw[->] (0,0) -- (0,\dzg@lt@yend);
  \ifdzg@lt@labels
    \node[right] at (\dzg@lt@xend,0) {$#1$};
    \node[above] at (0,\dzg@lt@yend) {$#3$};
    \foreach \dzg@x in {1,...,#2} {
      \draw (\dzg@x,0.12) -- (\dzg@x,-0.12);
      \node[below=2pt, font=\scriptsize] at (\dzg@x,0) {\dzg@x};
    }
    \foreach \dzg@y in {1,...,#4} {
      \draw (-0.12,\dzg@y) -- (0.12,\dzg@y);
      \node[left=4pt, font=\scriptsize] at (0,\dzg@y) {\dzg@y};
    }
  \fi
  \foreach \dzg@x/\dzg@y/\dzg@m in {#5}
    \node[vertex=\dzg@m] (-\dzg@x-\dzg@y) at (\dzg@x,\dzg@y) {};
}
